\chapter{Происхождение симметрии пространственного основного состояния}
\label{ch:v8-baryon-spatial-ground-state-symmetry-origin-gate}

Предыдущий гейт показал, что симметричная спин-ароматная ветвь требует
симметричного пространственного основного состояния. Проверим, достаточно
ли для этого одной перестановочной инвариантности.

\section{Полный перестановочно-инвариантный класс-суммовой тест}

Пусть $U_\pi$ --- представление перестановок трёх координат. Введём две
центральные суммы
\begin{equation}
 T_2=U_{(12)}+U_{(13)}+U_{(23)},\qquad
 T_3=U_{(123)}+U_{(132)}.
 \label{eq:v8-baryon-spatial-s3-class-sums}
\end{equation}
Каждый оператор
\begin{equation}
 H_{\alpha,\beta}=\alpha T_2+\beta T_3
 \label{eq:v8-baryon-spatial-central-hamiltonian}
\end{equation}
коммутирует со всем $S_3$. На тривиальном, знаковом и стандартном типах
его уровни равны
\begin{equation}
 E_{\mathbf1}=3\alpha+2\beta,\qquad
 E_{\mathbf1_{\rm sgn}}=-3\alpha+2\beta,\qquad
 E_{\mathbf2}=-\beta.
 \label{eq:v8-baryon-spatial-s3-levels}
\end{equation}
Поэтому все три типа могут быть основными:
\begin{equation}
 \begin{array}{c|ccc|c}
 (\alpha,\beta)&E_{\mathbf1}&E_{\mathbf1_{\rm sgn}}&E_{\mathbf2}
 &\text{основной тип}\\ \hline
 (-1,0)&-3&3&0&\mathbf1\\
 (1,0)&3&-3&0&\mathbf1_{\rm sgn}\\
 (0,1)&2&2&-1&\mathbf2
 \end{array}.
 \label{eq:v8-baryon-spatial-ground-branch-countermodels}
\end{equation}

\begin{theorem}[Недостаточность перестановочной инвариантности]
Условие $[H_{\rm space},U_\pi]=0$ не выбирает симметричное основное
состояние. Оно только запрещает смешивание неприводимых типов.
\end{theorem}

\section{Достаточное усиление}

Для скалярного координатного гамильтониана достаточным является условие,
что тепловая полугруппа улучшает положительность:
\begin{equation}
 f\ge0,\ f\ne0
 \quad\Longrightarrow\quad
 e^{-tH_{\rm space}}f>0
 \quad(t>0).
 \label{eq:v8-baryon-spatial-positivity-improving-condition}
\end{equation}
Тогда основное состояние единственно и может быть выбрано строго
положительным. Перестановочная инвариантность переводит его в другое
положительное основное состояние, а единственность вынуждает
\begin{equation}
 U_\pi\psi_0=\psi_0
 \qquad\text{для всех }\pi\in S_3.
 \label{eq:v8-baryon-spatial-positive-ground-symmetry}
\end{equation}
Это условно замыкает цепь: цветовой эпсилон плюс принцип Паули тогда
выбирают тривиальный спин-ароматный тип и магнитный множитель $-1/3$.

Однако текущий проект не содержит трёхкоординатного скалярного
$H_{\rm space}$, области его определения, кинетического члена и
потенциала, для которых можно было бы доказать улучшение положительности.
Ранее найденные полевые и шумовые операторы действуют на других носителях
и не могут быть переименованы в этот гамильтониан.

\begin{nogocriterion}[Граница пространственного селектора]
Перестановочная симметрия без невырожденности и улучшения положительности
не выбирает пространственный тип. Следовательно, симметричная
спин-ароматная ветвь остаётся условной до появления отдельного
координатного родителя.
\end{nogocriterion}

\begin{protocol}
\begin{enumerate}
 \item вычислительное поле: \textbf{$\mathbb Q(\alpha,\beta)$};
 \item класс-суммы $S_3$: \textbf{построены точно};
 \item перестановочная инвариантность: \textbf{сохранена во всех моделях};
 \item симметричная основная ветвь: \textbf{не вынуждена};
 \item знаковая и стандартная контрветви: \textbf{построены};
 \item достаточное условие: \textbf{улучшение положительности};
 \item координатный родитель проекта: \textbf{отсутствует};
 \item безусловный магнитный селектор: \textbf{не получен};
 \item следующий гейт: \textbf{итоговая электромагнитная проверка}.
\end{enumerate}
\end{protocol}

Машинный сертификат сохранён в
\path{s2t_v8_baryon_spatial_ground_state_symmetry_origin_gate_results.json}.